\chapter*{Resumen}
\addcontentsline{toc}{chapter}{Resumen}

\begingroup\singlespacing
El presente plan propone la \parteProceso{caracterización} mediante una \parteTecnologia{arquitectura analítica basada en clústeres espaciales} de los \parteContexto{incidentes de seguridad ciudadana en San Sebastián, Cusco}, registrados en SIPCOP Municipal. La arquitectura organiza la validación y preparación de los datos, el agrupamiento espacial, la evaluación estadística, el análisis temporal y la generación de resultados cartográficos. Se especifican sus componentes, interfaces, flujo de datos, reglas de trazabilidad y criterios de aceptación; una prueba de concepto reproducible permite evaluar su funcionamiento con el extracto disponible, que contiene 996 registros entre el 1 de enero y el 29 de agosto de 2025. El diagnóstico preliminar muestra coordenadas con resolución aproximada de 110 m y 412 ubicaciones distintas. Para el análisis se propone DBSCAN, comparado con HDBSCAN, y contrastes de agrupamiento frente a modelos nulos espaciales. La prueba de concepto será evaluada mediante casos funcionales, repetición controlada del flujo, consistencia de resultados y cumplimiento de requisitos. Los hallazgos del caso describen únicamente los registros disponibles; estos incluyen atenciones que no necesariamente constituyen delitos y no cubren la totalidad de incidentes del distrito.

\medskip
\noindent\textbf{Palabras clave:} ocurrencias de seguridad ciudadana, agrupamiento espacial, DBSCAN, \textit{hotspots}, patrullaje municipal, San Sebastián.
\endgroup

\clearpage
\chapter*{Abstract}
\addcontentsline{toc}{chapter}{Abstract}

\begingroup\singlespacing
This proposal addresses the \parteProceso{characterization}, through a \parteTecnologia{spatial-cluster-based analytical architecture}, of \parteContexto{public-safety incidents in San Sebastián, Cusco}, recorded in SIPCOP Municipal. The architecture organizes data validation and preparation, spatial clustering, statistical assessment, temporal analysis, and cartographic output. Its components, interfaces, data flow, traceability rules, and acceptance criteria are specified; a reproducible proof of concept evaluates its operation using the available extract of 996 records dated from 1 January to 29 August 2025. Preliminary diagnostics show coordinates with an approximate resolution of 110 m and 412 distinct locations. DBSCAN, compared with HDBSCAN, and spatial null-model tests are proposed for the case analysis. The proof of concept is evaluated through functional test cases, controlled workflow reruns, output consistency, and requirements compliance. Findings describe only the available records, which include attended events that are not necessarily crimes and do not represent all incidents in the district.

\medskip
\noindent\textbf{Keywords:} public security incidents, spatial clustering, DBSCAN, hotspots, municipal patrol, San Sebastián.
\endgroup
